% Preamble de notre livre de Mathematiques Terminale
\usepackage[utf8]{inputenc}
\usepackage[T1]{fontenc}
\usepackage[french]{babel}

% Polices et typographie
\usepackage{lmodern}
\usepackage{microtype}

% Mathematiques
\usepackage{amsmath, amssymb, amsfonts, mathtools, amsthm}
\usepackage{mathrsfs} % pour \mathscr
\usepackage{stmaryrd} % pour les doubles crochets

% Mise en page
\usepackage[top=2.5cm, bottom=2.5cm, left=2.2cm, right=2.2cm, headheight=15pt]{geometry}
\usepackage{fancyhdr}
\usepackage{titlesec}
\usepackage{enumitem}

% Graphiques et Figures
\usepackage{tikz}
\usetikzlibrary{shapes, positioning, arrows.meta, calc, patterns}
\usepackage{pgfplots}
\pgfplotsset{compat=1.18}

% Boites colorées (tcolorbox) pour Definition, Theoreme, Propriete, Methode
\usepackage{xcolor}
\usepackage[most]{tcolorbox}

% Couleurs harmonieuses
\definecolor{defcolor}{RGB}{31, 78, 121}      % Bleu profond pour les definitions
\definecolor{defbg}{RGB}{240, 244, 248}
\definecolor{thcolor}{RGB}{46, 117, 89}      % Vert foret pour les theoremes
\definecolor{thbg}{RGB}{241, 247, 244}
\definecolor{propcolor}{RGB}{128, 48, 108}   % Mauve/Prune pour les proprietes
\definecolor{propbg}{RGB}{252, 243, 250}
\definecolor{methcolor}{RGB}{204, 85, 0}      % Orange rouille pour les methodes
\definecolor{methbg}{RGB}{255, 250, 245}
\definecolor{excolor}{RGB}{100, 100, 100}     % Gris pour les exemples
\definecolor{exbg}{RGB}{250, 250, 250}

% Definition des styles tcolorbox
\newtcolorbox[auto counter, number within=chapter]{definition}[1][]{
  colback=defbg,
  colframe=defcolor,
  fonttitle=\bfseries\sffamily,
  title={Définition \thetcbcounter\ #1},
  arc=3pt,
  boxrule=1.2pt,
  left=10pt, right=10pt, top=8pt, bottom=8pt,
  before skip=12pt, after skip=12pt,
  enhanced,
  shadow={1.5pt}{-1.5pt}{0pt}{black!10}
}

\newtcolorbox[auto counter, number within=chapter]{theoreme}[1][]{
  colback=thbg,
  colframe=thcolor,
  fonttitle=\bfseries\sffamily,
  title={Théorème \thetcbcounter\ #1},
  arc=3pt,
  boxrule=1.2pt,
  left=10pt, right=10pt, top=8pt, bottom=8pt,
  before skip=12pt, after skip=12pt,
  enhanced,
  shadow={1.5pt}{-1.5pt}{0pt}{black!10}
}

\newtcolorbox[auto counter, number within=chapter]{propriete}[1][]{
  colback=propbg,
  colframe=propcolor,
  fonttitle=\bfseries\sffamily,
  title={Propriété \thetcbcounter\ #1},
  arc=3pt,
  boxrule=1.2pt,
  left=10pt, right=10pt, top=8pt, bottom=8pt,
  before skip=12pt, after skip=12pt,
  enhanced,
  shadow={1.5pt}{-1.5pt}{0pt}{black!10}
}

\newtcolorbox[auto counter, number within=chapter]{methode}[1][]{
  colback=methbg,
  colframe=methcolor,
  fonttitle=\bfseries\sffamily,
  title={Méthode \thetcbcounter\ #1},
  arc=3pt,
  boxrule=1.2pt,
  left=10pt, right=10pt, top=8pt, bottom=8pt,
  before skip=12pt, after skip=12pt,
  enhanced,
  shadow={1.5pt}{-1.5pt}{0pt}{black!10}
}

\newtcolorbox[auto counter, number within=chapter]{exemple}[1][]{
  colback=exbg,
  colframe=excolor,
  fonttitle=\bfseries\sffamily,
  title={Exemple \thetcbcounter\ #1},
  arc=3pt,
  boxrule=0.8pt,
  left=10pt, right=10pt, top=8pt, bottom=8pt,
  before skip=12pt, after skip=12pt,
  frame style={dash pattern=on 3pt off 3pt}
}

% Style listings pour Python
\usepackage{listings}
\lstdefinestyle{PythonStyle}{
  language=Python,
  basicstyle=\ttfamily\small,
  keywordstyle=\color{blue}\bfseries,
  commentstyle=\color{green!50!black}\itshape,
  stringstyle=\color{red!70!black},
  numbers=left,
  numberstyle=\tiny\color{gray},
  stepnumber=1,
  numbersep=8pt,
  backgroundcolor=\color{gray!5},
  showspaces=false,
  showstringspaces=false,
  frame=single,
  rulecolor=\color{gray!30},
  tabsize=4,
  captionpos=b,
  breaklines=true,
  breakatwhitespace=true,
  literate={é}{{\'e}}1 {è}{{\`e}}1 {à}{{\`a}}1 {ù}{{\`u}}1 {ç}{{\c{c}}}1 {œ}{{\oe}}1 {Œ}{{\OE}}1 {æ}{{\ae}}1 {Æ}{{\AE}}1 {î}{{\^i}}1 {ï}{{\"i}}1 {û}{{\^u}}1 {ô}{{\^o}}1,
}

% Configuration des en-têtes et pieds de page
\pagestyle{fancy}
\fancyhf{}
\fancyhead[LE,RO]{\sffamily\thepage}
\fancyhead[LO]{\sffamily\rightmark}
\fancyhead[RE]{\sffamily\leftmark}
\renewcommand{\headrulewidth}{0.4pt}
\renewcommand{\chaptermark}[1]{\markboth{\thechapter.\ #1}{}}
\renewcommand{\sectionmark}[1]{\markright{\thesection.\ #1}{}}

% Raccourcis mathematiques utiles
\newcommand{\R}{\mathbb{R}}
\newcommand{\N}{\mathbb{N}}
\newcommand{\Z}{\mathbb{Z}}
\newcommand{\Cnum}{\mathbb{C}} % Eviter conflit avec \C
\newcommand{\Eesp}{\mathbb{E}}
\newcommand{\Pprob}{\mathbb{P}}
\newcommand{\dd}{\mathrm{d}}
\newcommand{\iud}{\mathrm{i}}
\renewcommand{\binom}[2]{{C_{#1}^{#2}}}

% Configuration titlesec pour des chapitres et sections modernes
\titleformat{\chapter}[display]
  {\normalfont\huge\bfseries\sffamily\color{defcolor}}
  {\chaptertitlename\ \thechapter}
  {20pt}
  {\Huge}
\titleformat{\section}
  {\normalfont\Large\bfseries\sffamily\color{defcolor}}
  {\thesection}
  {1em}
  {}
\titleformat{\subsection}
  {\normalfont\large\bfseries\sffamily\color{thcolor}}
  {\thesubsection}
  {1em}
  {}
